% !TEX root = ../../../main.tex
% 所属：第二部分：多变量分析学 / 第2章：多元函数微分学
% 节：2.4 隐函数
% 正文按原文件顺序保存；由 ../chapter2.tex 调用。
\section{隐函数}

\begin{theorem}[隐函数存在定理]
	设 $z=F(x,y)$ 满足下列条件:
	\begin{enumerate}
		\item $F(x_0,y_0)=0$;
		\item $F(x,y)$ 在闭矩形
		\[
		D=\{(x,y):|x-x_0|\le a,\ |y-y_0|\le b\}
		\]
		上连续可微;
		\item $F_y(x_0,y_0)\ne0$.
	\end{enumerate}
	则有:
	\begin{enumerate}
		\item[(i)] 在 $(x_0,y_0)$ 的某邻域内, 方程 $F(x,y)=0$
		唯一确定一个隐函数
		\[
		y=f(x),\qquad x\in(x_0-\rho,x_0+\rho),\qquad \rho>0,
		\]
		且 $F(x,f(x))=0$, $f(x_0)=y_0$;
		\item[(ii)] $y=f(x)$ 在 $(x_0-\rho,x_0+\rho)$ 上连续;
		\item[(iii)] $y=f(x)$ 在该区间上可导, 且
		\[
		\frac{\dd y}{\dd x}
		=
		-\frac{F_x(x,y)}{F_y(x,y)}.
		\]
	\end{enumerate}
\end{theorem}

\begin{proof}
	由 $F_y(x_0,y_0)\ne0$, 不妨设
	$F_y(x_0,y_0)>0$. 又因 $F_y(x,y)$ 在 $D$ 上连续,
	存在以 $(x_0,y_0)$ 为中心的闭矩形邻域
	\[
	\overline D=
	\{(x,y):|x-x_0|\le\alpha,\ |y-y_0|\le\beta\}
	\subset D,
	\qquad
	0<\alpha<a,\quad 0<\beta<b,
	\]
	使得在 $\overline D$ 上恒有 $F_y(x,y)>0$.

	先令 $x=x_0$, 在区间
	$[y_0-\beta,y_0+\beta]$ 上考虑一元函数
	\[
	\varphi(y)=F(x_0,y).
	\]
	因为
	\[
	\varphi'(y)=F_y(x_0,y)>0,
	\qquad
	\varphi(y_0)=F(x_0,y_0)=0,
	\]
	所以
	\[
	F(x_0,y_0+\beta)>0,
	\qquad
	F(x_0,y_0-\beta)<0.
	\]
	在直线 $y=y_0+\beta$ 上, 函数
	$F(x,y_0+\beta)$ 关于 $x$ 连续, 且在 $x=x_0$ 时为正,
	故存在 $0<\rho_1<\alpha$, 使得
	\[
	F(x,y_0+\beta)>0,
	\qquad
	x\in[x_0-\rho_1,x_0+\rho_1].
	\]
	同理存在 $0<\rho_2<\alpha$, 使得
	\[
	F(x,y_0-\beta)<0,
	\qquad
	x\in[x_0-\rho_2,x_0+\rho_2].
	\]
	取
	\[
	\rho=\min\{\rho_1,\rho_2\}.
	\]
	于是对每个 $x\in(x_0-\rho,x_0+\rho)$,
	函数 $F(x,y)$ 在 $[y_0-\beta,y_0+\beta]$ 上连续,
	并在两端异号. 由介值定理, 至少存在一点
	$y\in(y_0-\beta,y_0+\beta)$ 使 $F(x,y)=0$.
	又因 $F_y(x,y)>0$, $F(x,y)$ 关于 $y$ 严格递增,
	故该零点唯一. 由此唯一确定
	\[
	y=f(x),
	\qquad
	F(x,f(x))=0,
	\qquad
	f(x_0)=y_0.
	\]

	
	下证 $\overline y=f(\overline x)$ 的连续性.
	$\forall\,\varepsilon>0$，考虑 $\overline y$ 的 $\varepsilon$-邻域，
	$\varepsilon$ 可取充分小，使得
	$(\overline y-\varepsilon,\overline y+\varepsilon)\subset(y_0-\beta,y_0+\beta)$.
	
	由连续性，
	
	$F(\overline x,\overline y+\varepsilon)>0$，
	$\exists\,\delta_1>0$，当
	$x\in(\overline x-\delta_1,\overline x+\delta_1)$ 时，
	$F(x,\overline y+\varepsilon)>0$.
	
	$F(\overline x,\overline y-\varepsilon)<0$，
	$\exists\,\delta_2>0$，
	$x\in(\overline x-\delta_2,\overline x+\delta_2)$ 时，
	$F(x,\overline y-\varepsilon)<0$.
	
	$\delta=\min\{\delta_1,\delta_2\}$，
	$x_*\in(\overline x-\delta,\overline x+\delta)$，
	$F(x_*,\overline y+\varepsilon)>0$，
	$F(x_*,\overline y-\varepsilon)<0$.
	\[
	\Rightarrow\quad
	y_*=f(x_*)\in(\overline y-\varepsilon,\overline y+\varepsilon),
	\qquad
	(x_*,y_*)\text{ 是 }F(x_*,y)=0\text{ 的唯一解}
	\quad\Rightarrow\quad
	|f(x_*)-f(\overline x)|<\varepsilon.
	\]
	
	综上，$\forall\,\varepsilon>0$，$\exists\,\delta>0$，
	$\forall\,x\in(\overline x-\delta,\overline x+\delta)$，
	$|f(x)-f(\overline x)|<\varepsilon$.

	

	最后证明 $f$ 的可导性. 仍令
	$\overline y=f(\overline x)$, 并取
	$\overline x+\Delta x\in(x_0-\rho,x_0+\rho)$. 记
	\[
	\Delta y=f(\overline x+\Delta x)-\overline y.
	\]
	由 $f$ 连续, 当 $\Delta x\to0$ 时有 $\Delta y\to0$.
	利用二元函数的中值公式, 存在 $\theta\in(0,1)$, 使得
	\[
	\begin{aligned}
	0
	={}&F(\overline x+\Delta x,\overline y+\Delta y)-F(\overline x,\overline y)\\
	={}&F_x(\overline x+\theta\Delta x,\overline y+\theta\Delta y)\Delta x\\
	&+F_y(\overline x+\theta\Delta x,\overline y+\theta\Delta y)\Delta y.
	\end{aligned}
	\]
	由于邻域内 $F_y>0$, 当 $\Delta x\ne0$ 时,
	\[
	\frac{\Delta y}{\Delta x}
	=
	-\frac{
	F_x(\overline x+\theta\Delta x,\overline y+\theta\Delta y)}
	{F_y(\overline x+\theta\Delta x,\overline y+\theta\Delta y)}.
	\]
	令 $\Delta x\to0$, 由 $F_x,F_y$ 连续, 得
	\[
	f'(\overline x)
	=
	-\frac{F_x(\overline x,\overline y)}{F_y(\overline x,\overline y)}.
	\]
	由于 $\overline x$ 任意, 结论成立.
\end{proof}
\subsection*{多元隐函数存在定理}
\begin{theorem}[多元隐函数存在定理]
	设 $n+1$ 元函数
	\[
	F(x_1,x_2,\ldots,x_n,y)
	\]
	满足下列条件:
	\begin{enumerate}
		\item
		\[
		F(x_1^0,x_2^0,\ldots,x_n^0,y^0)=0;
		\]
		\item 在闭长方体
		\[
		D=
		\left\{
		(x_1,x_2,\ldots,x_n,y):
		|y-y^0|\le b,\ 
		|x_i-x_i^0|\le a_i,\ i=1,2,\ldots,n
		\right\}
		\]
		上, 函数 $F$ 连续, 且具有连续偏导数
		$F_y,F_{x_i}\ (i=1,2,\ldots,n)$;
		\item
		\[
		F_y(x_1^0,x_2^0,\ldots,x_n^0,y^0)\ne0.
		\]
	\end{enumerate}
	则:
	\begin{enumerate}
		\item[(i)] 在点 $(x_1^0,x_2^0,\ldots,x_n^0,y^0)$ 附近,
		方程
		\[
		F(x_1,x_2,\ldots,x_n,y)=0
		\]
		唯一确定隐函数
		\[
		y=f(x_1,x_2,\ldots,x_n),
		\qquad
		(x_1,x_2,\ldots,x_n)
		\in O\bigl((x_1^0,x_2^0,\ldots,x_n^0),\rho\bigr),
		\]
		并且
		\[
		F\bigl(x_1,x_2,\ldots,x_n,
		f(x_1,x_2,\ldots,x_n)\bigr)=0,
		\qquad
		f(x_1^0,x_2^0,\ldots,x_n^0)=y^0;
		\]
		\item[(ii)] 隐函数 $y=f(x_1,x_2,\ldots,x_n)$
		在 $O\bigl((x_1^0,x_2^0,\ldots,x_n^0),\rho\bigr)$ 上连续;
		\item[(iii)] 隐函数 $y=f(x_1,x_2,\ldots,x_n)$
		在该邻域上具有连续偏导数, 且
		\[
		\frac{\partial y}{\partial x_i}
		=
		-\frac{F_{x_i}(x_1,x_2,\ldots,x_n,y)}
		{F_y(x_1,x_2,\ldots,x_n,y)},
		\qquad i=1,2,\ldots,n.
		\]
	\end{enumerate}
\end{theorem}

对于 $F(\vect x,y)=0$, 可以对 $x_i$ 求导:
\[
\frac{\partial F}{\partial x_i}
+
\frac{\partial F}{\partial y}
\frac{\partial y}{\partial x_i}
=0,
\]
由此即得上面的偏导数公式.



\subsection*{多个方程确定的隐函数}
我们对要研究的形式有一下的要求
\begin{enumerate}
	\item 方程有多个，比如$m$个。
	\item 未知数有多个，也是$m$个，与方程的数量相同.
\end{enumerate}
我们以以下的简单形式为例
\[
\begin{cases}
F(x,y,u,v)=0,\\
G(x,y,u,v)=0.
\end{cases}
\]

\begin{theorem}
	设 $F(x,y,u,v)$ 与 $G(x,y,u,v)$ 满足:
	\begin{enumerate}
		\item
		\[
		F(x_0,y_0,u_0,v_0)
		=
		G(x_0,y_0,u_0,v_0)
		=0;
		\]
		\item 在
		\[
		D=
		\left\{
		(x,y,u,v):
		|x-x_0|\le a,\ |y-y_0|\le b,\
		|u-u_0|\le c,\ |v-v_0|\le d
		\right\}
		\]
		上, $F,G$ 连续可微;
		\item 在 $(x_0,y_0,u_0,v_0)$ 处
		\[
		\frac{\partial(F,G)}{\partial(u,v)}
		=
		\begin{vmatrix}
		F_u&F_v\\
		G_u&G_v
		\end{vmatrix}
		\ne0.
		\]
	\end{enumerate}
	则在 $(x_0,y_0,u_0,v_0)$ 的某邻域内, 方程组唯一确定隐函数
	\[
	\begin{cases}
	u=f(x,y),\\
	v=g(x,y),
	\end{cases}
	\qquad
	(x,y)\in O((x_0,y_0),\rho),
	\qquad \rho>0.
	\]
	并且 $f(x,y),g(x,y)$ 连续可微.
\end{theorem}

\begin{proof}
	由
	\[
	\begin{vmatrix}
	F_u&F_v\\
	G_u&G_v
	\end{vmatrix}
	\ne0
	\]
	可知 $F_u,F_v,G_u,G_v$ 中至少有一个不为零.
	不妨设
	\[
	F_u(x_0,y_0,u_0,v_0)\ne0.
	\]
	由前面的隐函数存在定理, 在
	$(x_0,y_0,u_0,v_0)$ 的某邻域内,
	方程 $F(x,y,u,v)=0$ 唯一确定一个 $u$ 关于 $x,y,v$ 的隐函数:
	\[
	u=\varphi(x,y,v).
	\]
	$\varphi$ 连续可微, 对 $v$ 求导有
	\[
	\frac{\partial F}{\partial v}
	+
	\frac{\partial F}{\partial u}
	\frac{\partial\varphi}{\partial v}
	=0,
	\]
	因而
	\[
	\frac{\partial\varphi}{\partial v}
	=
	-\frac{F_v}{F_u}.
	\]

	将 $u=\varphi(x,y,v)$ 代入 $G$, 令
	\[
	H(x,y,v)
	=
	G\bigl(x,y,\varphi(x,y,v),v\bigr).
	\]
	则
	\[
	\begin{aligned}
	\frac{\partial H}{\partial v}
	&=
	\frac{\partial G}{\partial u}
	\frac{\partial\varphi}{\partial v}
	+\frac{\partial G}{\partial v}\\
	&=
	G_u\left(-\frac{F_v}{F_u}\right)+G_v\\
	&=
	\frac{F_uG_v-F_vG_u}{F_u}
	\ne0.
	\end{aligned}
	\]
	因此方程 $H(x,y,v)=0$ 又唯一确定 $v$ 关于 $x,y$ 的隐函数
	\[
	v=g(x,y).
	\]
	代回即得
	\[
	u=\varphi(x,y,g(x,y))=f(x,y).
	\]
	由于 $\varphi,g$ 连续可微, $f,g$ 也连续可微.

	由
	\[
	\begin{cases}
	F(x,y,u,v)=0,\\
	G(x,y,u,v)=0
	\end{cases}
	\]
	分别对 $x$ 求导,
	\[
	\begin{cases}
	F_x+F_u\dfrac{\partial u}{\partial x}
	+F_v\dfrac{\partial v}{\partial x}=0,\\[6pt]
	G_x+G_u\dfrac{\partial u}{\partial x}
	+G_v\dfrac{\partial v}{\partial x}=0.
	\end{cases}
	\]
	写成矩阵形式,
	\[
	\begin{pmatrix}
	F_u&F_v\\
	G_u&G_v
	\end{pmatrix}
	\begin{pmatrix}
	\dfrac{\partial u}{\partial x}\\[5pt]
	\dfrac{\partial v}{\partial x}
	\end{pmatrix}
	=
	-
	\begin{pmatrix}
	F_x\\
	G_x
	\end{pmatrix}.
	\]
	由于系数矩阵的行列式不为零, 上述方程组有唯一解,
	\[
	\begin{pmatrix}
	\dfrac{\partial u}{\partial x}\\[5pt]
	\dfrac{\partial v}{\partial x}
	\end{pmatrix}
	=
	-
	\begin{pmatrix}
	F_u&F_v\\
	G_u&G_v
	\end{pmatrix}^{-1}
	\begin{pmatrix}
	F_x\\
	G_x
	\end{pmatrix}.
	\]
	对 $y$ 同理. 合并可得
	\[
	\begin{pmatrix}
	\dfrac{\partial u}{\partial x}&
	\dfrac{\partial u}{\partial y}\\[6pt]
	\dfrac{\partial v}{\partial x}&
	\dfrac{\partial v}{\partial y}
	\end{pmatrix}
	=
	-
	\begin{pmatrix}
	F_u&F_v\\
	G_u&G_v
	\end{pmatrix}^{-1}
	\begin{pmatrix}
	F_x&F_y\\
	G_x&G_y
	\end{pmatrix}.
	\]
\end{proof}

\begin{theorem}[一般隐函数组存在定理]
	设 $m$ 个 $n+m$ 元函数
	\[
	F_i(x_1,x_2,\ldots,x_n,y_1,y_2,\ldots,y_m),
	\qquad i=1,2,\ldots,m,
	\]
	满足以下条件:
	\begin{enumerate}
		\item
		\[
		F_i(x_1^0,x_2^0,\ldots,x_n^0,y_1^0,y_2^0,\ldots,y_m^0)
		=0,
		\qquad i=1,2,\ldots,m;
		\]
		\item 在闭长方体
		\[
		\begin{aligned}
		D=\bigl\{&(x_1,x_2,\ldots,x_n,y_1,y_2,\ldots,y_m):\\
		&|x_i-x_i^0|\le a_i,\ i=1,2,\ldots,n;\ 
		|y_j-y_j^0|\le b_j,\ j=1,2,\ldots,m
		\bigr\}
		\end{aligned}
		\]
		上, 函数 $F_i\ (i=1,2,\ldots,m)$ 连续,
		且具有连续偏导数;
		\item 在
		$(x_1^0,x_2^0,\ldots,x_n^0,y_1^0,y_2^0,\ldots,y_m^0)$
		处, Jacobi 行列式
		\[
		\frac{\partial(F_1,F_2,\ldots,F_m)}
		{\partial(y_1,y_2,\ldots,y_m)}
		\ne0.
		\]
	\end{enumerate}
	则:
	\begin{enumerate}
		\item[(i)] 在该点的某个邻域内, 方程组
		\[
		\begin{cases}
		F_1(x_1,\ldots,x_n,y_1,\ldots,y_m)=0,\\
		F_2(x_1,\ldots,x_n,y_1,\ldots,y_m)=0,\\
		\qquad\vdots\\
		F_m(x_1,\ldots,x_n,y_1,\ldots,y_m)=0
		\end{cases}
		\]
		唯一确定向量值隐函数
		\[
		\begin{pmatrix}
		y_1\\y_2\\ \vdots\\y_m
		\end{pmatrix}
		=
		\begin{pmatrix}
		f_1(x_1,x_2,\ldots,x_n)\\
		f_2(x_1,x_2,\ldots,x_n)\\
		\vdots\\
		f_m(x_1,x_2,\ldots,x_n)
		\end{pmatrix},
		\qquad
		(x_1,x_2,\ldots,x_n)
		\in O\bigl((x_1^0,x_2^0,\ldots,x_n^0),\rho\bigr),
		\]
		满足原方程组, 并且
		\[
		y_i^0=f_i(x_1^0,x_2^0,\ldots,x_n^0),
		\qquad i=1,2,\ldots,m;
		\]
		\item[(ii)] 该向量值隐函数在上述邻域上连续;
		\item[(iii)] 该向量值隐函数在上述邻域上具有连续偏导数,
\end{enumerate}
\end{theorem}
导数有
\[
\begin{pmatrix}
	\dfrac{\partial y_1}{\partial x_1}&
	\dfrac{\partial y_1}{\partial x_2}&\cdots&
	\dfrac{\partial y_1}{\partial x_n}\\
	\dfrac{\partial y_2}{\partial x_1}&
	\dfrac{\partial y_2}{\partial x_2}&\cdots&
	\dfrac{\partial y_2}{\partial x_n}\\
	\vdots&\vdots&\ddots&\vdots\\
	\dfrac{\partial y_m}{\partial x_1}&
	\dfrac{\partial y_m}{\partial x_2}&\cdots&
	\dfrac{\partial y_m}{\partial x_n}
\end{pmatrix}
=
-
\begin{pmatrix}
	\dfrac{\partial F_1}{\partial y_1}&
	\dfrac{\partial F_1}{\partial y_2}&\cdots&
	\dfrac{\partial F_1}{\partial y_m}\\
	\dfrac{\partial F_2}{\partial y_1}&
	\dfrac{\partial F_2}{\partial y_2}&\cdots&
	\dfrac{\partial F_2}{\partial y_m}\\
	\vdots&\vdots&\ddots&\vdots\\
	\dfrac{\partial F_m}{\partial y_1}&
	\dfrac{\partial F_m}{\partial y_2}&\cdots&
	\dfrac{\partial F_m}{\partial y_m}
\end{pmatrix}^{-1}
\begin{pmatrix}
	\dfrac{\partial F_1}{\partial x_1}&
	\dfrac{\partial F_1}{\partial x_2}&\cdots&
	\dfrac{\partial F_1}{\partial x_n}\\
	\dfrac{\partial F_2}{\partial x_1}&
	\dfrac{\partial F_2}{\partial x_2}&\cdots&
	\dfrac{\partial F_2}{\partial x_n}\\
	\vdots&\vdots&\ddots&\vdots\\
	\dfrac{\partial F_m}{\partial x_1}&
	\dfrac{\partial F_m}{\partial x_2}&\cdots&
	\dfrac{\partial F_m}{\partial x_n}
\end{pmatrix}.
\]
\subsection*{多元函数的反函数}
以$\R^2 \mapsto \R^2$为例：
设 $D\subset\mathbb{R}^2$ 为开集, $f:D\to\mathbb{R}^2$ 为映射,
\[
f(u,v)=\bigl(x(u,v),y(u,v)\bigr),
\qquad (u,v)\in D,
\]
其中 $x(u,v),y(u,v)$ 在 $D$ 内有连续偏导数. 若在
$(u_0,v_0)\in D$ 处, 对应点为
\[
x_0=x(u_0,v_0),\qquad y_0=y(u_0,v_0),
\]
且
\[
\frac{\partial(x,y)}{\partial(u,v)}\bigg|_{(u_0,v_0)}
=\begin{vmatrix}
\dfrac{\partial x}{\partial u}&\dfrac{\partial x}{\partial v}\\[4pt]
\dfrac{\partial y}{\partial u}&\dfrac{\partial y}{\partial v}
\end{vmatrix}_{(u_0,v_0)}\ne0,
\]
则在 $(u_0,v_0)$ 的某邻域内, 映射 $f$ 存在反函数
\[
u=u(x,y),\qquad v=v(x,y),
\]
且该反函数具有连续偏导数. 由复合函数求导法则,
\[
\begin{pmatrix}
\dfrac{\partial x}{\partial u}&\dfrac{\partial x}{\partial v}\\[5pt]
\dfrac{\partial y}{\partial u}&\dfrac{\partial y}{\partial v}
\end{pmatrix}
\begin{pmatrix}
\dfrac{\partial u}{\partial x}&\dfrac{\partial u}{\partial y}\\[5pt]
\dfrac{\partial v}{\partial x}&\dfrac{\partial v}{\partial y}
\end{pmatrix}
=\begin{pmatrix}1&0\\0&1\end{pmatrix},
\]
从而
\[
\begin{pmatrix}
\dfrac{\partial u}{\partial x}&\dfrac{\partial u}{\partial y}\\[5pt]
\dfrac{\partial v}{\partial x}&\dfrac{\partial v}{\partial y}
\end{pmatrix}
=\begin{pmatrix}
\dfrac{\partial x}{\partial u}&\dfrac{\partial x}{\partial v}\\[5pt]
\dfrac{\partial y}{\partial u}&\dfrac{\partial y}{\partial v}
\end{pmatrix}^{-1},
\qquad
\frac{\partial(u,v)}{\partial(x,y)}
=\left(\frac{\partial(x,y)}{\partial(u,v)}\right)^{-1}.
\]

